% GENERATED FILE. Edit canonical lessons, then run npm run generate:books.
% canonical-source-hash: fnv1a64:98b3245f2a138d34
% canonical-lessons: KA-C207-musu, KA-C207-netuhaku, KA-C207-kuttu, KA-C207-perisu, KA-C207-sangrahisu

\chapter{To sniff, To hang up, To pound, To stack, To collect}
\label{ch:ka-to-sniff-to-hang-up-to-pound-to-stack-to-collect}
\hlchaptermodality{207}

\begin{chapteropening}
\textbf{By the end of this chapter:} \emph{I can say to sniff, to hang up, to pound, to stack and to collect.}

Complete the last lesson of chapter 207: I can say to sniff, to hang up, to pound, to stack and to collect.
\end{chapteropening}

\section[\texorpdfstring{\kn{ಮೂಸು} (mūsu)}{ಮೂಸು (mūsu)}]{\kn{ಮೂಸು} (mūsu) — to sniff}

\label{lesson:KA-C207-musu}



\begin{quote}
Before the new one: say the Kannada for to escape, then the Kannada for to hurry.
\end{quote}

\subsection*{You'll want to know: \kn{ಮೂಸು}}
\textbf{\kn{ಮೂಸು}} — \emph{mūsu} — \textquotedblleft{}to sniff\textquotedblright{}.

\begin{grammarlens}[title={What you've built}]
One more word for this chapter.
\end{grammarlens}

\subsection*{Your turn}
\begin{itemize}
\raggedright
  \item \emph{Say it:} \emph{mūsu}
  \item \emph{Say it:} \emph{mūsu}, once more
  \item \emph{Say it:} say \emph{mūsu}
  \item \emph{Recall:} say the Kannada for to escape, then the Kannada for to hurry, then say \emph{mūsu} again
\end{itemize}

\begin{tcolorbox}[breakable,colback=teal!4,colframe=teal!35!black,title={Before you move on}]
What does \kn{ಮೂಸು} mean? (\textquotedblleft{}To sniff\textquotedblright{}.) Say it once more.
\end{tcolorbox}
\section[\texorpdfstring{\kn{ನೇತುಹಾಕು} (nētuhāku)}{ನೇತುಹಾಕು (nētuhāku)}]{\kn{ನೇತುಹಾಕು} (nētuhāku) — to hang up}

\label{lesson:KA-C207-netuhaku}



\begin{quote}
Before the new one: say the Kannada for to hurry, then the Kannada for to sniff.
\end{quote}

\subsection*{You'll want to know: \kn{ನೇತುಹಾಕು}}
\textbf{\kn{ನೇತುಹಾಕು}} — \emph{nētuhāku} — \textquotedblleft{}to hang up\textquotedblright{}.

\begin{grammarlens}[title={What you've built}]
One more word for this chapter.
\end{grammarlens}

\subsection*{Your turn}
\begin{itemize}
\raggedright
  \item \emph{Say it:} \emph{nētuhāku}
  \item \emph{Say it:} \emph{nētuhāku}, once more
  \item \emph{Say it:} say \emph{nētuhāku}
  \item \emph{Recall:} say the Kannada for to hurry, then the Kannada for to sniff, then say \emph{nētuhāku} again
\end{itemize}

\begin{tcolorbox}[breakable,colback=teal!4,colframe=teal!35!black,title={Before you move on}]
What does \kn{ನೇತುಹಾಕು} mean? (\textquotedblleft{}To hang up\textquotedblright{}.) Say it once more.
\end{tcolorbox}
\section[\texorpdfstring{\kn{ಕುಟ್ಟು} (kuṭṭu)}{ಕುಟ್ಟು (kuṭṭu)}]{\kn{ಕುಟ್ಟು} (kuṭṭu) — to pound (in a mortar)}

\label{lesson:KA-C207-kuttu}



\begin{quote}
Before the new one: say the Kannada for to sniff, then the Kannada for to hang up.
\end{quote}

\subsection*{You'll want to know: \kn{ಕುಟ್ಟು}}
\textbf{\kn{ಕುಟ್ಟು}} — \emph{kuṭṭu} — \textquotedblleft{}to pound (in a mortar)\textquotedblright{}.

\begin{grammarlens}[title={What you've built}]
One more word for this chapter.
\end{grammarlens}

\subsection*{Your turn}
\begin{itemize}
\raggedright
  \item \emph{Say it:} \emph{kuṭṭu}
  \item \emph{Say it:} \emph{kuṭṭu}, once more
  \item \emph{Say it:} say \emph{kuṭṭu}
  \item \emph{Recall:} say the Kannada for to sniff, then the Kannada for to hang up, then say \emph{kuṭṭu} again
\end{itemize}

\begin{tcolorbox}[breakable,colback=teal!4,colframe=teal!35!black,title={Before you move on}]
What does \kn{ಕುಟ್ಟು} mean? (\textquotedblleft{}To pound (in a mortar)\textquotedblright{}.) Say it once more.
\end{tcolorbox}
\section[\texorpdfstring{\kn{ಪೇರಿಸು} (pērisu)}{ಪೇರಿಸು (pērisu)}]{\kn{ಪೇರಿಸು} (pērisu) — to stack, to pile up}

\label{lesson:KA-C207-perisu}



\begin{quote}
Before the new one: say the Kannada for to hang up, then the Kannada for to pound.
\end{quote}

\subsection*{You'll want to know: \kn{ಪೇರಿಸು}}
\textbf{\kn{ಪೇರಿಸು}} — \emph{pērisu} — \textquotedblleft{}to stack, to pile up\textquotedblright{}.

\begin{grammarlens}[title={What you've built}]
One more word for this chapter.
\end{grammarlens}

\subsection*{Your turn}
\begin{itemize}
\raggedright
  \item \emph{Say it:} \emph{pērisu}
  \item \emph{Say it:} \emph{pērisu}, once more
  \item \emph{Say it:} say \emph{pērisu}
  \item \emph{Recall:} say the Kannada for to hang up, then the Kannada for to pound, then say \emph{pērisu} again
\end{itemize}

\begin{tcolorbox}[breakable,colback=teal!4,colframe=teal!35!black,title={Before you move on}]
What does \kn{ಪೇರಿಸು} mean? (\textquotedblleft{}To stack, to pile up\textquotedblright{}.) Say it once more.
\end{tcolorbox}
\section[\texorpdfstring{\kn{ಸಂಗ್ರಹಿಸು} (saṅgrahisu)}{ಸಂಗ್ರಹಿಸು (saṅgrahisu)}]{\kn{ಸಂಗ್ರಹಿಸು} (saṅgrahisu) — to collect, to gather}

\label{lesson:KA-C207-sangrahisu}



\begin{quote}
Before the new one: say the Kannada for to pound, then the Kannada for to stack.
\end{quote}

\subsection*{You'll want to know: \kn{ಸಂಗ್ರಹಿಸು}}
\textbf{\kn{ಸಂಗ್ರಹಿಸು}} — \emph{saṅgrahisu} — \textquotedblleft{}to collect, to gather\textquotedblright{}.

\begin{grammarlens}[title={What you've built}]
One more word for this chapter.
\end{grammarlens}

\subsection*{Your turn}
\begin{itemize}
\raggedright
  \item \emph{Say it:} \emph{saṅgrahisu}
  \item \emph{Say it:} \emph{saṅgrahisu}, once more
  \item \emph{Say it:} say \emph{saṅgrahisu}
  \item \emph{Recall:} say the Kannada for to pound, then the Kannada for to stack, then say \emph{saṅgrahisu} again
\end{itemize}

\begin{tcolorbox}[breakable,colback=teal!4,colframe=teal!35!black,title={Before you move on}]
What does \kn{ಸಂಗ್ರಹಿಸು} mean? (\textquotedblleft{}To collect, to gather\textquotedblright{}.) Say it once more.
\end{tcolorbox}
